\documentclass[11pt]{article}
\usepackage[a4paper,margin=1in]{geometry}
\usepackage{amsmath,amssymb,amsthm,microtype}
\newtheorem{theorem}{Theorem}[section]
\theoremstyle{remark}\newtheorem{remark}[theorem]{Remark}
\title{Annular \(H^\infty\) and Hardy Criteria for the Direct Weighted Prefix}
\author{Bin Wang}\date{July 2026}
\begin{document}\maketitle
\begin{abstract}
The remaining determinant bridge is an absolute weighted coefficient norm,
but it need not be proved coefficient by coefficient or root by root.  We
give two direct analytic criteria.  If the logarithmic mismatch is small in
$H^\infty$ on any circle strictly outside the target disk, Cauchy's estimate
controls the full weighted coefficient norm.  A vanishing $H^2$ boundary
norm on the same annulus also suffices by Cauchy--Schwarz.  Both bounds are
explicit.  At radius $1.41$, which lies between the target radius $1.4$ and
the certified deterministic radius $1.426787\ldots$, their constants are
$139.007\ldots$ and $8.292\ldots$.  An endpoint polynomial family shows that
vanishing $H^2$ norm on the target circle alone is insufficient.  No actual
noisy annular convergence is asserted.
\end{abstract}

\section{Logarithmic mismatch}
Let
\[
 g_\sigma(z)=\sum_{n\ge2}
 \frac{\tau_{\sigma,n}-a_n}{n}z^n
\]
be holomorphic on a neighborhood of the closed disk $|z|\le\rho$, where
$\rho>R>0$.  Define the full absolute coefficient norm
\[
 P_\sigma^\infty(R)=\sum_{n\ge2}
 \frac{|\tau_{\sigma,n}-a_n|R^n}{n}.
\]

\section{Two annular criteria}
\begin{theorem}\label{thm:analytic}
Put $x=R/\rho<1$.
\begin{enumerate}
\item If $\sup_{|z|\le\rho}|g_\sigma(z)|\le M_\sigma$, then
\[
 P_\sigma^\infty(R)\le M_\sigma\frac{x^2}{1-x}.
\]
\item If
\[
 \|g_\sigma\|_{H^2(\rho)}^2:=
 \sum_{n\ge2}
 \frac{|\tau_{\sigma,n}-a_n|^2\rho^{2n}}{n^2}
 \le H_\sigma^2,
\]
then
\[
 P_\sigma^\infty(R)\le
 H_\sigma\frac{x^2}{\sqrt{1-x^2}}.
\]
\end{enumerate}
In particular, $M_\sigma\to0$ or $H_\sigma\to0$ makes the full norm vanish.
Every moving prefix is bounded by $P_\sigma^\infty(R)$, so the RH-292
shortened prefix follows.
\end{theorem}
\begin{proof}
Cauchy's coefficient estimate gives
\[
 \frac{|\tau_{\sigma,n}-a_n|}{n}\le M_\sigma\rho^{-n},
\]
and summing the geometric series proves the first bound.  For the second,
apply Cauchy--Schwarz to
\[
 \sum_{n\ge2}
 \left(
 \frac{|\tau_{\sigma,n}-a_n|\rho^n}{n}
 \right)x^n.
\]
The square sum of $x^n$ from $n=2$ onward is $x^4/(1-x^2)$.
\end{proof}

\section{Endpoint obstruction}
\begin{theorem}\label{thm:endpoint}
The strict radius gap in the $H^2$ criterion cannot be removed uniformly.
For every $N\ge1$, let
\[
 g_N(z)=\frac1N\sum_{n=2}^{N+1}(z/R)^n.
\]
Then
\[
 \|g_N\|_{H^2(R)}=N^{-1/2}\longrightarrow0,
\qquad
 \sum_{n=2}^{N+1}|[z^n]g_N|R^n=1.
\]
\end{theorem}
\begin{proof}
There are $N$ coefficients, each with boundary-scaled modulus $1/N$.
Their square sum is $1/N$, while their absolute sum is one.
\end{proof}

\section{Certified annulus and protocol}
The deterministic coefficient radius is
$\rho_*=0.85\lambda=1.426787483864073\ldots$.  Taking $\rho=1.41$ and
$R=1.4$ gives
\[
 \frac{x^2}{1-x}=139.0070921986\ldots,\qquad
 \frac{x^2}{\sqrt{1-x^2}}=8.2924678943\ldots.
\]
The computation records these constants and the endpoint family.  It does
not estimate $g_\sigma$ for the noisy operator.

\begin{remark}
The theorem is a reopening criterion, not a verified annular bridge.  If
activated it supplies the direct $P_\sigma^\infty$ bound, but no such
$H^\infty$ or $H^2$ convergence is currently proved.  Gates A--E remain
false/open and no
Hilbert--Polya, zeta-divisor, or RH statement follows.
\end{remark}
\end{document}
